%% Custom Pandoc LaTeX template for "Building an LLM from Zero"
%% Minimal dependencies — works with texlive-xetex + texlive-latex-extra
\documentclass[11pt,oneside]{book}

%% Encoding & fonts
\usepackage{fontspec}
\usepackage{unicode-math}
\defaultfontfeatures{Scale=MatchLowercase}

%% Page geometry
\usepackage[margin=1in,top=1.1in,bottom=1.1in]{geometry}

%% Colors
\usepackage[dvipsnames,svgnames,x11names]{xcolor}
\definecolor{linkblue}{HTML}{006699}
\definecolor{tealcolor}{HTML}{0097A7}
\definecolor{codebg}{HTML}{F5F5F5}

%% Hyperlinks + embedded PDF metadata (read by Google Scholar, Zotero, Mendeley)
\usepackage[
  colorlinks=true,
  linkcolor=linkblue,
  urlcolor=linkblue,
  citecolor=linkblue,
  bookmarks=true,
  bookmarksnumbered=true
]{hyperref}

%% Embed document metadata into the PDF file itself.
%% This is what Google Scholar and citation managers read directly
%% from the PDF — separate from the visual title page.
\hypersetup{
  pdftitle={$title$$if(subtitle)$: $subtitle$$endif$},
  pdfauthor={$author$},
  pdfsubject={A hands-on tutorial for building a GPT-style language model from scratch using Python and PyTorch. No GPU required.},
  pdfkeywords={large language model, GPT, transformer, PyTorch, deep learning, machine learning, NLP, tutorial, artificial intelligence},
  pdfcreator={Pandoc + XeLaTeX},
  pdfproducer={jackluu.io}
}

%% Code highlighting (Pandoc)
$if(highlighting-macros)$
$highlighting-macros$
$endif$

%% Inline code styling
\usepackage{fancyvrb}
\usepackage{listings}

%% Better tables
\usepackage{longtable}
\usepackage{booktabs}
\usepackage{array}

%% Graphics
\usepackage{graphicx}

%% List spacing
\usepackage{enumitem}
\setlist{noitemsep, topsep=4pt}

%% Pandoc tight list macro (required — Pandoc injects \tightlist for compact lists)
\providecommand{\tightlist}{%
  \setlength{\itemsep}{0pt}\setlength{\parskip}{0pt}}

%% Part numbering: use Roman numerals (I, II, III)
\renewcommand{\thepart}{\Roman{part}}

%% Section heading styles
\usepackage{titlesec}

%% Part page style
\titleformat{\part}[display]
  {\centering\normalfont\Huge\bfseries\color{tealcolor}}
  {Part \thepart}{20pt}{\Huge}

%% Chapter style: "Chapter 1" label above, then title
\titleformat{\chapter}[display]
  {\normalfont\huge\bfseries\color{tealcolor}}
  {Chapter \thechapter}{16pt}{\Huge}
\titlespacing*{\chapter}{0pt}{20pt}{20pt}

\titleformat{\section}
  {\normalfont\Large\bfseries\color{tealcolor!80!black}}
  {\thesection}{1em}{}
\titleformat{\subsection}
  {\normalfont\large\bfseries}
  {\thesubsection}{1em}{}

%% Header / footer
\usepackage{fancyhdr}
\pagestyle{fancy}
\fancyhf{}
\fancyhead[LE,RO]{\small\textit{Building an LLM from Zero}}
\fancyhead[RE,LO]{\small\textit{Truong (Jack) Luu}}
\fancyfoot[C]{\small\thepage}
\renewcommand{\headrulewidth}{0.4pt}

%% Paragraph spacing
\setlength{\parindent}{0pt}
\setlength{\parskip}{8pt}

%% Blockquote styling
\usepackage{mdframed}
\newmdenv[
  topline=false, bottomline=false, rightline=false,
  linewidth=3pt, linecolor=tealcolor,
  backgroundcolor=tealcolor!5,
  skipabove=8pt, skipbelow=8pt,
  innerleftmargin=10pt, innerrightmargin=8pt
]{blockquotebox}
\renewenvironment{quote}{\begin{blockquotebox}\itshape}{\end{blockquotebox}}

%% Table of contents depth
\setcounter{tocdepth}{2}

%% ---------------------------------------------------------------
%% DOCUMENT
%% ---------------------------------------------------------------
\begin{document}

%% Title page
\begin{titlepage}
  \newgeometry{margin=1.5in}
  \vspace*{2cm}
  {\color{tealcolor}\rule{\linewidth}{3pt}}
  \vspace{1cm}

  {\fontsize{32}{38}\selectfont\bfseries $title$\par}
  \vspace{0.5cm}

  $if(subtitle)$
  {\fontsize{16}{20}\selectfont\color{gray} $subtitle$\par}
  \vspace{1cm}
  $endif$

  {\color{tealcolor}\rule{\linewidth}{1pt}}
  \vspace{1.5cm}

  {\large $author$\par}
  \vspace{0.3cm}
  $if(first-draft)$
  {\color{gray}\small First draft: $first-draft$\par}
  \vspace{0.15cm}
  $endif$
  {\color{gray}\small Last edited: $date$\par}

  \vfill
  {\small\color{gray}
    Source code: \href{https://github.com/jackluucoding/build-llm-from-zero}{github.com/jackluucoding/build-llm-from-zero}\\
    Read online: \href{https://jackluu.io/book/}{jackluu.io/book/}\\
    Contact: \href{mailto:hi@jackluu.io}{hi@jackluu.io}\\[6pt]
    \textcopyright\ 2025 Truong (Jack) Luu.
    Code: MIT License. Book text: CC BY-NC 4.0.
  }
  \restoregeometry
\end{titlepage}

%% Abstract page — helps Google Scholar identify and index this work
\clearpage
\thispagestyle{empty}
\vspace*{3cm}
{\large\bfseries Abstract}
\vspace{0.5em}

\noindent
This book is a hands-on tutorial for building a GPT-style language model entirely from scratch
using Python and PyTorch. It covers every component of the Transformer architecture —
tokenization, embeddings, self-attention, multi-head attention, feed-forward layers, residual
connections, and layer normalization — before assembling them into a complete language model,
training it on the Tiny Shakespeare dataset, and generating text. No GPU is required; the full
training run completes in approximately 20-30 minutes on a standard CPU-only laptop. The book
is designed for readers with basic Python knowledge and no prior machine learning background,
including high school students, college instructors seeking classroom-ready materials, IT
professionals, and parents teaching their children about artificial intelligence.

\vspace{1em}
\noindent\textbf{Keywords:} large language model, GPT, Transformer, PyTorch, deep learning,
natural language processing, machine learning tutorial, artificial intelligence education

\clearpage

%% Table of contents
\tableofcontents
\clearpage

%% \mainmatter resets page numbering to arabic and starts chapter counter from 1
\mainmatter

%% Main content
$body$

\end{document}
